\begin{helper}
Let $n\ge1$ be an integer, let $1<p<\infty$ and let
$A\colon\mathbb{Z}\to\mathbb{C}$.  Write $T_k=T_k(n,p,A,\cdot)$ for the $k$-th
summand of the series (19) as in \noderef{Thm2Term}, $\lambda_k=\lambda_k(n,p)$
as in \noderef{LambdaCoef} and $\epsilon_k=\epsilon_k(n,p)$ as in
\noderef{EpsCoef}.

For a word $w=(v_1,\dots,v_q)$ of complex letters define the iterated
directional derivative $D_w$ on functions $f\colon\mathbb{C}\to\mathbb{C}$ by
\[
  D_{()}f=f,\qquad
  D_{(v,v_1,\dots,v_q)}f=\bigl(z\mapsto \mathrm{d}\bigl(D_{(v_1,\dots,v_q)}f\bigr)(z)[v]\bigr),
\]
so that each letter contributes one real Fr\'echet differentiation in its own
direction, the leading letter being the outermost one; and put
\[
  V_w=\prod_{i=1}^{q}|v_i| .
\]
Then such a pair $(D,V)$ exists, satisfies the two recursions above together
with $V_{()}=1$, $V_{(v,v_1,\dots,v_q)}=|v|\,V_{(v_1,\dots,v_q)}$ and
$V_w\ge0$, and has the following two properties.  For every word $w$ of length
$q=|w|$, every integer $m\ge0$ and every $\xi\ne0$,
\[
  \bigl|D_wT_{n+m}(\xi)\bigr|
  \le\Bigl(2\bigl((n+m)\bigl(3+\sqrt{p-1}\bigr)+q\bigr)\Bigr)^{q}
      \,V_w\,\bigl(2|A_{n+m}|\bigr)\,
      |\xi|^{\lambda_{n+m}-q},
\]
and $D_wT_{n+m}$ is of class $C^\infty$ at every $\xi\ne0$.

This is the paper's "we note that $\lim\lambda_k=\infty$.  Therefore, we can
perform any finite number of term-by-term differentiations in formula (19)"
(\texttt{paper.tex:387}) made quantitative on a single summand.  Each
differentiation lowers the common exponent of the summand by exactly one and
inflates its coefficient by a factor at most
$2\bigl((n+m)(3+\sqrt{p-1})+q\bigr)$, which is polynomial in $n+m$; against
the geometric decay of $\rho^{\lambda_{n+m}}$ supplied by the linear lower
bound for $\lambda_k$, such a factor is harmless at every order.  The
operators $D$ and $V$ are produced existentially rather than as separate
definition nodes because only the stated recursions are ever used.
\end{helper}
\begin{proof}
Throughout, $n\ge1$ and $1<p<\infty$ are fixed and $k=n+m$.  All assertions
about $\xi$ are for $\xi\ne0$, where $|\xi|^{a}$ is an ordinary real power of
a positive real and $\xi^{j}$, $\bar\xi^{\,l}$ are ordinary integer powers of a
nonzero complex number; in particular all the functions below are defined and
$C^\infty$ on the open set $\mathbb{C}-\{0\}$.

\emph{Step 0: the operators, and the estimates for $\lambda_k$ and
$\epsilon_k$.}  Define
\[
  D_wf=\text{the right fold of }w\text{ under }(v,g)\mapsto
    \bigl(z\mapsto\mathrm{d}g(z)[v]\bigr)\text{ starting from }f,
  \qquad V_w=\prod_{i}|v_i| ,
\]
which satisfy $D_{()}f=f$ and $D_{(v)\frown w}f=(z\mapsto\mathrm{d}(D_wf)(z)[v])$
by definition of a fold, and likewise $V_{()}=1$ and
$V_{(v)\frown w}=|v|\,V_w$; and $V_w\ge0$ because it is a product of moduli,
by induction on $w$.

From \noderef{LambdaEpsEstimates} we record, for every $m\ge0$,
\begin{equation}\tag{1}
  \lambda_{n+m}\ge0,\qquad
  |\epsilon_{n+m}|\le1,\qquad
  \lambda_{n+m}\le\bigl((n:\mathbb{R})+m\bigr)\sqrt{p-1}.
\end{equation}
For $m=0$ these are (23): $\lambda_n=0$ and $\epsilon_n=0$, and the third
reads $0\le n\sqrt{p-1}$, which holds since $n\ge0$ and $p-1>0$.  For $m\ge1$
we have $n+1\le n+m$, so (24) gives $|\epsilon_{n+m}|<1$ and (25) gives both
$\lambda_{n+m}\ge(p-1)/(np)>0$ and $\lambda_{n+m}\le(n+m)\sqrt{p-1}$.

\emph{Step 1: the monomial class and its first derivatives.}  For $s\in
\mathbb{R}$ and $j,l\in\mathbb{Z}$ put
\[
  E_{s,j,l}(\xi)=|\xi|^{s}\,\xi^{j}\,\bar\xi^{\,l}\qquad(\xi\ne0).
\]
Each factor is $C^\infty$ on $\mathbb{C}-\{0\}$ --- $|\xi|^{s}$ because
$\xi\mapsto|\xi|$ is $C^\infty$ and positive there and $t\mapsto t^{s}$ is
$C^\infty$ on $(0,\infty)$; $\xi^{j}$ and $\bar\xi^{\,l}$ because they are
quotients of polynomials in $\xi$, $\bar\xi$ with nonvanishing denominators
--- so $E_{s,j,l}$ is $C^\infty$ on $\mathbb{C}-\{0\}$.  Its Wirtinger
derivatives are
\begin{equation}\tag{2}
  \partial_\xi E_{s,j,l}=\Bigl(\frac{s}{2}+j\Bigr)E_{s,j-1,l},
  \qquad
  \partial_{\bar\xi}E_{s,j,l}=\Bigl(\frac{s}{2}+l\Bigr)E_{s,j,l-1} .
\end{equation}
Indeed $|\xi|^{s}=(\xi\bar\xi)^{s/2}$, so by
\noderef{WirtingerModulusPowerRule}
$\partial_\xi|\xi|^{s}=\tfrac{s}{2}|\xi|^{s}\xi^{-1}$ and
$\partial_{\bar\xi}|\xi|^{s}=\tfrac{s}{2}|\xi|^{s}\bar\xi^{\,-1}$; by
\noderef{WirtingerPowerRule} and \noderef{WirtingerConjRule}
$\partial_\xi\xi^{j}=j\xi^{j-1}$, $\partial_{\bar\xi}\xi^{j}=0$,
$\partial_\xi\bar\xi^{\,l}=0$ and
$\partial_{\bar\xi}\bar\xi^{\,l}=l\bar\xi^{\,l-1}$.  Applying
\noderef{WirtingerProductRule} twice and collecting terms gives (2), using
$\xi^{j}\xi^{-1}=\xi^{j-1}$ and $\bar\xi^{\,l}\bar\xi^{\,-1}=\bar\xi^{\,l-1}$.

\emph{Step 2: finite combinations, and the three bookkeeping quantities.}  By
a \emph{configuration} we mean a finite list
$L=\bigl((c_1,j_1,l_1),\dots,(c_N,j_N,l_N)\bigr)$ with $c_i\in\mathbb{C}$ and
$j_i,l_i\in\mathbb{Z}$.  For $s\in\mathbb{R}$ put
\[
  S_{s,L}(\xi)=\sum_{i=1}^{N}c_i\,E_{s,j_i,l_i}(\xi),
  \qquad
  \|L\|=\sum_{i=1}^{N}|c_i| .
\]
Say that $L$ \emph{has exponent total $e$} (relative to $s$) if
$s+j_i+l_i=e$ for every $i$, and that $L$ \emph{has frequency bound $K$} if
$\tfrac12|s|+|j_i|+|l_i|\le K$ for every $i$.  Two elementary consequences:

(2a) if $L$ has exponent total $e$ then
\begin{equation}\tag{3}
  \bigl|S_{s,L}(\xi)\bigr|\le\|L\|\,|\xi|^{e}\qquad(\xi\ne0),
\end{equation}
because $|E_{s,j_i,l_i}(\xi)|=|\xi|^{s}|\xi|^{j_i}|\xi|^{l_i}
=|\xi|^{s+j_i+l_i}=|\xi|^{e}$ and the triangle inequality applies to the
finite sum;

(2b) $S_{s,L}$ is $C^\infty$ on $\mathbb{C}-\{0\}$, being a finite linear
combination of the $E_{s,j,l}$ of Step 1.

\emph{Step 3: one differentiation of a configuration.}  Fix $s$, a letter
$v\in\mathbb{C}$ and a configuration $L$, and define a new configuration of
length $2N$,
\[
  \mathcal{D}_{s,v}L
  =\Bigl(\bigl(c_i(\tfrac{s}{2}+j_i)v,\;j_i-1,\;l_i\bigr)\Bigr)_{i\le N}
   \frown
   \Bigl(\bigl(c_i(\tfrac{s}{2}+l_i)\bar v,\;j_i,\;l_i-1\bigr)\Bigr)_{i\le N}.
\]
For a function $F$ that is real-differentiable at $\xi$ one has
$\mathrm{d}F(\xi)[v]=v\,\partial_\xi F(\xi)+\bar v\,\partial_{\bar\xi}F(\xi)$,
which is just the definition of the Wirtinger derivatives
$\partial_\xi=\tfrac12(\partial_1-i\partial_2)$,
$\partial_{\bar\xi}=\tfrac12(\partial_1+i\partial_2)$ of \noderef{DZ} and
\noderef{DZbar} read backwards, together with the $\mathbb{R}$-linearity of
$v\mapsto\mathrm{d}F(\xi)[v]$ and $v=v_1+iv_2$.  Applying this to
$F=S_{s,L}$, which is differentiable at every $\xi\ne0$ by (2b), and using
(2) termwise,
\begin{equation}\tag{4}
  \mathrm{d}S_{s,L}(\xi)[v]=S_{s,\mathcal{D}_{s,v}L}(\xi)\qquad(\xi\ne0).
\end{equation}
Moreover:
\begin{enumerate}
\item[(i)] if $L$ has exponent total $e$ then $\mathcal{D}_{s,v}L$ has
  exponent total $e-1$: the two families of exponent pairs are
  $(j_i-1,l_i)$ and $(j_i,l_i-1)$, and in either case the sum
  $s+j+l$ drops by exactly $1$;
\item[(ii)] if $L$ has frequency bound $K$ then $\mathcal{D}_{s,v}L$ has
  frequency bound $K+1$, since $|j_i-1|\le|j_i|+1$ and $|l_i-1|\le|l_i|+1$;
\item[(iii)] if $L$ has frequency bound $K$ with $K\ge0$ then
  \begin{equation}\tag{5}
    \bigl\|\mathcal{D}_{s,v}L\bigr\|\le 2K\,|v|\,\|L\| ,
  \end{equation}
  because $\bigl|c_i(\tfrac{s}{2}+j_i)v\bigr|=|c_i|\,
  \bigl|\tfrac{s}{2}+j_i\bigr|\,|v|\le|c_i|\,K\,|v|$ by
  $\bigl|\tfrac{s}{2}+j_i\bigr|\le\tfrac12|s|+|j_i|\le K$, and likewise
  $\bigl|c_i(\tfrac{s}{2}+l_i)\bar v\bigr|\le|c_i|\,K\,|v|$ using
  $|\bar v|=|v|$; summing the two families gives $2K|v|\|L\|$.
\end{enumerate}

\emph{Step 4: the induction on the length of the word.}  We claim: for every
word $w$, every $s,e,K\in\mathbb{R}$ with $K\ge0$ and every configuration $L$
with exponent total $e$ and frequency bound $K$, there is a configuration $L'$
with
\begin{equation}\tag{6}
  \begin{aligned}
  &D_w\bigl(S_{s,L}\bigr)(\xi)=S_{s,L'}(\xi)\quad(\xi\ne0),\\
  &L'\text{ has exponent total }e-|w|\text{ and frequency bound }K+|w|,\\
  &\|L'\|\le\bigl(2(K+|w|)\bigr)^{|w|}\,V_w\,\|L\| .
  \end{aligned}
\end{equation}
We induct on $w$.  If $w$ is empty, take $L'=L$: then $D_{()}S_{s,L}=S_{s,L}$,
the exponent total and frequency bound are $e-0$ and $K+0$, and
$\|L\|\le1\cdot1\cdot\|L\|$ since $(2K)^{0}=1$ and $V_{()}=1$.

Let $w=(v)\frown w_0$ and let $L_0$ be a configuration provided by the
inductive hypothesis for $w_0$, so that $D_{w_0}S_{s,L}=S_{s,L_0}$ off the
origin, $L_0$ has exponent total $e-|w_0|$ and frequency bound $K+|w_0|$, and
$\|L_0\|\le(2(K+|w_0|))^{|w_0|}V_{w_0}\|L\|$.  Put
$L'=\mathcal{D}_{s,v}L_0$.  Fix $\xi\ne0$.  The two functions
$D_{w_0}S_{s,L}$ and $S_{s,L_0}$ agree on the open set $\mathbb{C}-\{0\}$,
which is a neighbourhood of $\xi$, so they have the same Fr\'echet derivative
at $\xi$; hence, by the defining recursion of $D$ and by (4),
\[
  D_wS_{s,L}(\xi)=\mathrm{d}\bigl(D_{w_0}S_{s,L}\bigr)(\xi)[v]
  =\mathrm{d}S_{s,L_0}(\xi)[v]=S_{s,\mathcal{D}_{s,v}L_0}(\xi)=S_{s,L'}(\xi).
\]
By (i) and (ii) of Step 3, $L'$ has exponent total
$(e-|w_0|)-1=e-|w|$ and frequency bound $(K+|w_0|)+1=K+|w|$.  Finally, by (5)
applied with the frequency bound $K+|w_0|\ge0$ of $L_0$, and then by the
inductive bound on $\|L_0\|$,
\[
  \|L'\|\le2\bigl(K+|w_0|\bigr)|v|\,\|L_0\|
  \le2\bigl(K+|w_0|\bigr)|v|\,
     \bigl(2(K+|w_0|)\bigr)^{|w_0|}V_{w_0}\|L\|
  =\bigl(2(K+|w_0|)\bigr)^{|w|}V_{w}\|L\| ,
\]
using $|v|\,V_{w_0}=V_w$ and $|w|=|w_0|+1$.  Since
$0\le2(K+|w_0|)\le2(K+|w|)$, raising to the $|w|$-th power preserves the
inequality, and $V_w\|L\|\ge0$, so
$\|L'\|\le(2(K+|w|))^{|w|}V_w\|L\|$.  This proves (6).

\emph{Step 5: the summand as a configuration.}  Fix $m\ge0$ and put $k=n+m$,
\[
  s=\lambda_k+n-k=\lambda_{n+m}-m,\qquad
  K=(n+m)\bigl(3+\sqrt{p-1}\bigr),\qquad
  L_k=\bigl((A_k,\;m,\;0),\;(\epsilon_k\overline{A_k},\;-n,\;k)\bigr).
\]
By \noderef{Thm2Term}, for $\xi\ne0$,
\[
  T_k(\xi)=\bigl(A_k\xi^{k}+\epsilon_k\overline{A_k}\bar\xi^{\,k}\bigr)
    |\xi|^{s}\xi^{-n}
  =A_k|\xi|^{s}\xi^{k-n}
   +\epsilon_k\overline{A_k}\,|\xi|^{s}\xi^{-n}\bar\xi^{\,k},
\]
and $k-n=m$, so
\begin{equation}\tag{7}
  T_k(\xi)=S_{s,L_k}(\xi)\qquad(\xi\ne0).
\end{equation}
We check the three bookkeeping quantities of $L_k$.

\emph{Exponent total $\lambda_k$.}  For the first entry
$s+m+0=\lambda_k+n-k+m=\lambda_k$ since $k=n+m$; for the second
$s+(-n)+k=\lambda_k+n-k-n+k=\lambda_k$.

\emph{Frequency bound $K$.}  Write $N=(n:\mathbb{R})+m$, so $K=3N+N\sqrt{p-1}$
and, by (1), $0\le\lambda_k\le N\sqrt{p-1}$.  Since $\lambda_k\ge0$ and
$m\ge0$ we get $-( \lambda_k+m)\le s=\lambda_k-m\le\lambda_k+m$, i.e.
\begin{equation}\tag{8}
  |s|\le\lambda_k+m\le N\sqrt{p-1}+m .
\end{equation}
For the first entry, $\tfrac12|s|+|m|+|0|\le\tfrac12 N\sqrt{p-1}
+\tfrac12m+m\le N\sqrt{p-1}+3N=K$, because $\tfrac12m+m=\tfrac32m\le3N$ (as
$m\le N$) and $\tfrac12N\sqrt{p-1}\le N\sqrt{p-1}$.  For the second entry,
$\tfrac12|s|+|-n|+|k|\le\tfrac12N\sqrt{p-1}+\tfrac12m+n+N$, and
$\tfrac12m+n+N\le\tfrac12N+N+N\le3N$ using $m\le N$ and $n\le N$; so again the
total is at most $N\sqrt{p-1}+3N=K$.  Note $K\ge0$.

\emph{Mass at most $2|A_k|$.}  By (1), $|\epsilon_k|\le1$, so
\[
  \|L_k\|=|A_k|+\bigl|\epsilon_k\overline{A_k}\bigr|
  =|A_k|+|\epsilon_k|\,|A_k|\le|A_k|+|A_k|=2|A_k| .
\]

\emph{Step 6: conclusion.}  Fix a word $w$ with $|w|=q$, an integer $m\ge0$
and $\xi\ne0$, and keep $k$, $s$, $K$, $L_k$ as in Step 5.  The two functions
$T_k$ and $S_{s,L_k}$ agree on the open set $\mathbb{C}-\{0\}$ by (7), so by
the same local-agreement argument used in Step 4 (applied $q$ times, i.e. by
the evident induction on $w$) one has $D_wT_k=D_wS_{s,L_k}$ at every point of
$\mathbb{C}-\{0\}$.  Let $L'$ be the configuration produced by Step 4 for the
word $w$ and the data $(s,\lambda_k,K,L_k)$, legitimate because $K\ge0$ and
$L_k$ has exponent total $\lambda_k$ and frequency bound $K$.  Then
$D_wT_k(\xi)=S_{s,L'}(\xi)$, $L'$ has exponent total $\lambda_k-q$ and
\[
  \|L'\|\le\bigl(2(K+q)\bigr)^{q}V_w\|L_k\|
  \le\bigl(2(K+q)\bigr)^{q}V_w\cdot2|A_k| ,
\]
the last step because $(2(K+q))^{q}V_w\ge0$.  Hence by (3),
\[
  \bigl|D_wT_k(\xi)\bigr|=\bigl|S_{s,L'}(\xi)\bigr|
  \le\|L'\|\,|\xi|^{\lambda_k-q}
  \le\Bigl(2\bigl((n+m)(3+\sqrt{p-1})+q\bigr)\Bigr)^{q}V_w
     \bigl(2|A_{n+m}|\bigr)|\xi|^{\lambda_{n+m}-q},
\]
which is the asserted bound.

For the smoothness, $S_{s,L'}$ is $C^\infty$ on $\mathbb{C}-\{0\}$ by (2b), in
particular $C^\infty$ at $\xi$; and $D_wT_k$ agrees with $S_{s,L'}$ on the
open set $\mathbb{C}-\{0\}$, which is a neighbourhood of $\xi$.  Being
$C^\infty$ at a point depends only on the germ of the function there, so
$D_wT_{n+m}$ is $C^\infty$ at $\xi$.  As $w$, $m$ and $\xi\ne0$ were
arbitrary, both conclusions hold, and the pair $(D,V)$ constructed in Step 0
witnesses the existential statement.
\end{proof}
